\documentclass[10pt,a4paper]{amsart}
\usepackage[margin=19mm]{geometry}
\usepackage{amsmath,amssymb,mathtools}
\usepackage[scr=rsfs,cal=euler]{mathalfa}
\usepackage{xfrac}
\usepackage[unicode,hidelinks,bookmarks=false]{hyperref}
\newcommand{\ie}{i.e.\ }
\newcommand{\cf}{cf.\ }
\newcommand{\resp}{respectively\ }
\newcommand{\Subsection}{\subsection}
\providecommand{\nobreakdash}{\nobreak\hbox{-}}
\setlength{\emergencystretch}{2em}
\multlinegap0pt
\usepackage{bm}
\usepackage{bbm}
\usepackage{dsfont}
\usepackage{xspace}
\usepackage{cancel}
\usepackage{mathtools}
\usepackage{amsthm}
\makeatletter
\@ifundefined{proposition}{\newtheorem{proposition}{Proposition}}{}
\makeatother
\title[Exterior Dirichlet Problems for Hessian Quotient Equations o]{Exterior Dirichlet Problems for Hessian Quotient Equations of Mixed Type}
\author{Yu Lei, Zhisu Li}
\date{Source: arXiv:2608.10834v1 (CC BY 4.0)}
\begin{document}
\maketitle

\noindent\emph{Source and licence.} Exterior Dirichlet Problems for Hessian Quotient Equations of Mixed Type. Version arXiv:2608.10834v1 (CC BY 4.0), distributed under the Creative Commons Attribution 4.0 International licence, as recorded in the arXiv licence metadata for this exact version and shown on the abstract page https://arxiv.org/abs/2608.10834v1. Full rights notice: licenses/arxiv-2608.10834-CC-BY-4.0.txt.\par\smallskip

\noindent\textit{Editorial scope.} Counted unit: the Proposition whose source label is \texttt{None} in the article (source0.tex lines 1083--1088, source label \texttt{None}), delivered together with the article's own complete original proof of it (source0.tex lines 1090--1461, one \texttt{proof} environment of 372 lines). No mathematics is written, repaired or re-derived: statement and proof are the article's own text line for line. Required context retained uncounted: prop (lines 320--354). Every definition, display and label of those lines is retained unchanged. Omitted from this package and not redistributed: 1--319, 355--1082, 1089--1089, 1462--1639. Nothing inside a retained excerpt is omitted or altered: every retained body is delivered exactly as the article's lines are written, so no omission receipt of any kind accompanies this package. Citations: the retained excerpts carry no citation command at all, so no bibliography entry of the article is transcribed and no reference is redistributed. Reference adaptation: none was needed and none was made; every reference inside the delivered unit resolves inside it, so no unresolved reference or citation is produced. Printed slips kept literal: no printed word, symbol, hypothesis or number of the counted statement or of its proof was altered. Layout and class: the article's own class is \texttt{amsart} with packages including amsmath, amstext, amssymb, amsfonts, amscd, amsthm, geometry, esint, hyperref, mathdots, mathrsfs, enumerate, subfigure, extarrows; the wrapper is the lane's pinned \texttt{amsart} editorial wrapper at 10pt on A4, which restates the theorem-like environments the retained text needs and adds the article's own macro definitions that the retained text needs (none), with extra wrapper packages (bm, bbm, dsfont, xspace, cancel, mathtools). The delivered numbering is the unit's own. Assets: no retained span contains a figure environment, a graphics-inclusion command, a PDF-page inclusion, a table or a TikZ picture; no image file is copied into this package, the package declares no asset dependency and no image file is redistributed with it. Licence: version arXiv:2608.10834v1 of this article is distributed under the Creative Commons Attribution 4.0 International licence, as recorded in the arXiv licence metadata for this exact version and shown on the abstract page https://arxiv.org/abs/2608.10834v1. A full rights notice is delivered at licenses/arxiv-2608.10834-CC-BY-4.0.txt. Characters: every retained excerpt is pure ASCII, so the delivered engine, XeLaTeX with its default Latin Modern text font, reports no missing character.
	\begin{proposition}\label{prop}
		The following properties hold:
		\begin{enumerate}
			\item $\widetilde{\Gamma}_k$ are convex cones and
			\[
			\Gamma_1 = \widetilde{\Gamma}_1 \supset \widetilde{\Gamma}_2 \supset \cdots \supset \widetilde{\Gamma}_n \supset \Gamma_2.
			\]
			
			\item If $\lambda = (\lambda_1, \ldots, \lambda_n) \in \widetilde{\Gamma}_k$, then
			\[
			\frac{\partial \left[ \frac{\sigma_k(\eta)}{\sigma_l(\eta)} \right]}{\partial \lambda_i}
			= \sum_{p \neq i} \frac{\sigma_{k-1}(\eta|p)\sigma_l(\eta) - \sigma_k(\eta)\sigma_{l-1}(\eta|p)}{\sigma_l(\eta)^2},
			\]
			and
			\[
			\frac{\partial \left[ \frac{\sigma_k(\eta)}{\sigma_l(\eta)} \right]}{\partial \lambda_i}
			\geq \frac{n(k-l)}{k(n-l)} \sum_{p \neq i} \frac{\sigma_{k-1}(\eta|p)\sigma_l(\eta|p)}{\sigma_l(\eta)^2},
			\]
			for any $i = 1, 2, \ldots, n$, where $0 \leq l < k \leq n$.
			
			\item If $\lambda = (\lambda_1, \ldots, \lambda_n) \in \widetilde{\Gamma}_k$, then $\left[ \frac{\sigma_k(\eta)}{\sigma_l(\eta)} \right]^{\frac{1}{k-l}}$ $(0 \leq l < k \leq n)$ is concave with respect to $\lambda$. Hence, for any $(\xi_1, \ldots, \xi_n)$, we have
			\[
			\sum_{i,j} \frac{\partial^2 \left[ \frac{\sigma_k(\eta)}{\sigma_l(\eta)} \right]}{\partial \lambda_i \partial \lambda_j} \xi_i \xi_j
			\leq
			\left(1 - \frac{1}{k-l}\right)
			\frac{
				\left(
				\sum_i \frac{\partial \left[ \frac{\sigma_k(\eta)}{\sigma_l(\eta)} \right]}{\partial \lambda_i} \xi_i
				\right)^2
			}{
				\frac{\sigma_k(\eta)}{\sigma_l(\eta)}
			}.
			\]
		\end{enumerate}
	\end{proposition}

	\begin{proposition}
		There exists $C > 0$, independent of $R$, such that
		\[
		|D^2 u_R| \le C \quad \text{on } \partial \Omega_R.
		\]
	\end{proposition}

	\begin{proof}
		We estimate the second derivatives on $\partial D$ and on $\partial E_R$ separately.
		
		\paragraph{Inner boundary $\partial D$.}
		Fix $p \in \partial D$. Choose local coordinates $y = (y', y_n)$ centered at $p$ such that
		\[
		(\mathbb{R}^n \setminus \overline{D}) \cap B_r(0) = \{ y_n > \rho(y') \} \cap B_r(0), \qquad \rho(0) = 0, \quad D\rho(0) = 0.
		\]
		Write $u := u_R$. Since $u = \Phi$ on $\partial D$,
		\[
		u(y', \rho(y')) = \Phi(y', \rho(y')) \quad \text{for } |y'| \ll 1.
		\]
		Differentiating once gives
		\[
		u_\alpha + u_n \rho_\alpha = \Phi_\alpha + \Phi_n \rho_\alpha.
		\]
		At $y=0$, because $D\rho(0)=0$, we get $u_\alpha(0) = \Phi_\alpha(0)$. Differentiating again yields
		\[
		u_{\alpha\beta}(0) = \Phi_{\alpha\beta}(0) - u_n(0)\rho_{\alpha\beta}(0).
		\]
		Since $|u_n| \le C$, we obtain $|u_{\alpha\beta}(0)| \le C$.
		
		
		For the mixed derivatives.
		Define the linearized operator
		\begin{equation}
			\mathcal{L} := F^{ij}(D^2 u_R)\partial_{ij}. \label{eq:L_mixed}
		\end{equation}
		Since $F(D^2 u_R)=1$ is constant, differentiating the equation gives
		\begin{equation}
			\mathcal{L}(\partial_\alpha u_R)=0 \qquad (\alpha=1,\dots,n). \label{eq:Lder_mixed}
		\end{equation}
		In local coordinates at $p$, we have
		$$D_\delta := \{ x \in \mathbb{R}^n \setminus \overline{D} : |x| < \delta \}.$$
		We first prove a key estimate for a barrier function. 
		Since $w_R$ is a strict subsolution and $u_R$ is a solution of the approximating problem, 
		the comparison principle gives $u_R-w_R>0$ in $\Omega_R$, and Hopf's lemma yields 
		$\partial_\nu(u_R-w_R)>0$ on $\partial D$. 
		We have
		\begin{equation}
			0<F(D^2w_R)-F(D^2u_R)
			\le F^{ij}(D^2u_R)(w_R-u_R)_{ij}
			= \mathcal{L}(w_R-u_R). \label{eq:concavity_zeta_mixed}
		\end{equation}
		Since $\mathcal{L}u_R=0$, this gives $\mathcal{L}w_R>0$, hence
		\begin{equation}
			\mathcal{L}(u_R-w_R)=-\mathcal{L}w_R<0 \quad \text{in } D_\delta. \label{eq:Luw_zeta_mixed}
		\end{equation}
		
		Now set
		\[
		\zeta:=u_R-w_R+td-\frac{N}{2}d^2.
		\]
		Using \eqref{eq:Luw_zeta_mixed}, the smoothness of $d$, 
		$$
		\begin{aligned}
			\mathcal L\left(\frac{d^2}{2}\right)
			&= F^{ij}\partial_i d\,\partial_j d + d\,\mathcal L d \\
			&\ge \lambda_0\sum_i F^{ii} - d\,C_d\sum_i F^{ii} \\
			&= (\lambda_0 - C C_d)\sum_i F^{ii}.
		\end{aligned}$$
		Since $0<d<\delta$ in $D_\delta$, taking $\delta>0$ sufficiently small so that $\lambda_0 - C_d\delta \ge \lambda_0/2$, we obtain
		\begin{equation}
			\mathcal L\left(\frac{d^2}{2}\right) \ge \frac{\lambda_0}{2}\sum_i F^{ii}.
		\end{equation}
		we obtain, for $t>0$ sufficiently small and $N>0$ sufficiently large,
		\begin{equation}
			\begin{aligned}
				\mathcal{L}\zeta
				&= \mathcal{L}(u_R-w_R) + t\mathcal{L}d - N\mathcal{L}\left(\frac{d^2}{2}\right) \\
				&\le tC_d\sum_i F^{ii} - \frac{N\lambda_0}{2}\sum_i F^{ii} \\
				&= \left(tC_d-\frac{N\lambda_0}{2}\right)\sum_i F^{ii} \\
				&\le  -\sum_iF^{ii},
			\end{aligned}
			\label{eq:Lzeta_mixed}
		\end{equation}
		for $tC_d-\frac{N\lambda_0}{2}\leq -1$ and $C_d = n\max_{\partial D}|\kappa_\alpha|$. 
		Moreover, by Hopf's lemma, $u_R-w_R\ge cd$ near $\partial D$. Taking $\delta>0$ sufficiently small, we have $\zeta\ge0$ on $\partial D_\delta$.
		
		Now define the auxiliary function
		\[
		\Psi := A_1\zeta + A_2 d - A_3\sum_{\beta=1}^{n-1}|\partial_\beta(u_R-w_R)|^2,
		\]
		where $A_1,A_2,A_3>0$ are constants to be chosen with $A_1\gg A_2\gg A_3\gg 1$.
		
		We estimate $\mathcal{L}\Psi$ term by term. From \eqref{eq:Lzeta_mixed},
		\[
		\mathcal{L}(A_1\zeta) \le -A_1\sum_i F^{ii}.
		\]
		Since $d$ is smooth, there exists $C_d>0$ such that
		\[
		\mathcal{L}(A_2d) \le A_2 C_d \sum_i F^{ii}.
		\]
		For the last term, set
		\[
		v_\beta:=\partial_\beta(u_R-w_R) \qquad (\beta=1,\dots,n-1).
		\]
		From \eqref{eq:Lder_mixed}, $\mathcal{L}v_\beta = -\mathcal{L}(\partial_\beta w_R)$. 
		Since $w_R=u_{\mathrm{near}}$ near $\partial D$ is fixed and smooth, there exists $C_2>0$ such that
		\[
		|\mathcal{L}v_\beta| \le C_2 \sum_i F^{ii}.
		\]
		By the $C^1$ estimate, $|v_\beta|\le C$, hence
		\[
		|2v_\beta \mathcal{L}v_\beta| \le C_3 \sum_i F^{ii}.
		\]
		By the ellipticity estimate in Proposition \ref{prop}, there exists $c_0>0$ such that
		\[
		\sum_{\beta<n} F^{ij}\partial_i v_\beta \partial_j v_\beta \ge c_0 \sum_i F^{ii}.
		\]
		Therefore,
		\[
		\mathcal{L}\left(\sum_{\beta<n}v_\beta^2\right)
		= 2\sum_{\beta<n}F^{ij}\partial_i v_\beta \partial_j v_\beta
		+ 2\sum_{\beta<n}v_\beta \mathcal{L}v_\beta
		\ge (2c_0 - C_3)\sum_i F^{ii}.
		\]
		Consequently,
		\[
		\mathcal{L}\left(-A_3\sum_{\beta<n}v_\beta^2\right)
		\le -A_3(2c_0 - C_3)\sum_i F^{ii}.
		\]
		Combining the above estimates, we obtain
		\begin{align}
			\mathcal{L}\Psi
			&\le -A_1 \sum_iF^{ii}
			+ A_2C_d\sum_i F^{ii}
			- A_3(2c_0 - C_3)\sum_i F^{ii} \notag \\
			&= -\left(A_1- A_2C_d + A_3(2c_0 - C_3)\right)\sum_i F^{ii}. \label{eq:LPsi_est_mixed}
		\end{align}
		Choosing $A_1\gg A_2+A_3$ so that
		\[
		A_1 - A_2C_d + A_3(2c_0 - C_3) \ge \frac{A_1}{2},
		\]
		we get
		\begin{equation}
			\mathcal{L}\Psi \le -\frac{A_1}{2}\sum_i F^{ii} < 0 \quad \text{in } D_\delta. \label{eq:LPsi_mixed}
		\end{equation}
		On $\partial D\cap B_\delta$, we have $d=0$, $u_R=w_R$, and $\partial_\beta(u_R-w_R)=0$, so $\Psi=0$. 
		On $\partial B_\delta\cap D_\delta$, by Hopf's lemma, $u_R-w_R\ge cd$ near $\partial D$. Hence
		\[
		\zeta = u_R-w_R+td-\frac{N}{2}d^2 \ge cd - \frac{N}{2}d^2 \ge d\left(c-\frac{N}{2}\delta\right)>0
		\]
		by taking $\delta>0$ sufficiently small. Thus $\Psi=A_1\zeta+A_2d-A_3\sum|\partial_\beta v|^2\ge0$ for $A_1$ sufficiently large. 
		Therefore $\Psi\ge0$ on $\partial D_\delta$. By the maximum principle, together with \eqref{eq:LPsi_mixed},
		\begin{equation}
			\Psi \ge 0 \quad \text{in } D_\delta. \label{eq:Psi_bound_mixed}
		\end{equation}
		Now fix $\alpha\in\{1,\dots,n-1\}$. For either choice of sign $\pm$, define
		\[
		\Psi_\pm := \Psi \pm \partial_\alpha(u_R-w_R).
		\]
		We verify that $\Psi_\pm$ satisfies the hypotheses of the maximum principle.
		Indeed, by \eqref{eq:Lder_mixed} and \eqref{eq:LPsi_mixed},
		\[
		\mathcal{L}(\Psi_\pm)
		= \mathcal{L}\Psi \pm \mathcal{L}(\partial_\alpha u_R) \mp \mathcal{L}(\partial_\alpha w_R)
		= \mathcal{L}\Psi \mp \mathcal{L}(\partial_\alpha w_R) < 0 \quad \text{in } D_\delta,
		\]
		after possibly increasing $A_1$ to absorb the term $\mathcal{L}(\partial_\alpha w_R)$.
		On the boundary $\partial D_\delta$, by \eqref{eq:Psi_bound_mixed} and the $C^1$ bound,
		\[
		\Psi_\pm = \Psi \pm \partial_\alpha(u_R-w_R) \ge 0 \quad \text{on } \partial D_\delta.
		\]
		Therefore, by the maximum principle,
		\begin{equation}
			\Psi \pm \partial_\alpha(u_R-w_R) \ge 0 \quad \text{in } D_\delta. \label{eq:Psi_alpha_bound_mixed}
		\end{equation}
		At $p\in\partial D$, we have $\Psi(0)=0$ and $\partial_\alpha(u_R-w_R)(0)=0$. 
		Applying Hopf's lemma to $\Psi_+$ and $\Psi_-$ yields
		\[
		\partial_n\Psi(0) \pm \partial_n\partial_\alpha(u_R-w_R)(0) \ge 0.
		\]
		Thus
		\[
		|u_{\alpha n}(0) - w_{\alpha n}(0)| \le \partial_n\Psi(0).
		\]
		Since $w_R=u_{\mathrm{near}}$ near $\partial D$ is fixed and smooth, and $\Psi$ is composed of $C^1$-bounded functions, we have
		\[
		|u_{\alpha n}(0)| \le |w_{\alpha n}(0)| + |\partial_n\Psi(0)| \le C,
		\]
		where $C$ is independent of $R$. Since $p\in\partial D$ and $\alpha\in\{1,\dots,n-1\}$ are arbitrary,
		\begin{equation}
			|u_{\alpha n}| \le C \quad \text{on } \partial D. \label{eq:mixed_mixed}
		\end{equation}
		
		
		
		
		We now estimate $u_{nn}(0)$ for $k\leq n-1$. By the boundedness of tangential and mixed derivatives,
		\[
		\Delta u \, I_n - D^2u = 
		\begin{pmatrix}
			\Delta u - u_{11} & -u_{12} & \cdots & -u_{1n} \\
			-u_{21} & \Delta u - u_{22} & \cdots & -u_{2n} \\
			\vdots & \vdots & \ddots & \vdots \\
			-u_{n1} & -u_{n2} & \cdots & \Delta u - u_{nn}
		\end{pmatrix}.
		\]
		It follows that 
		\[
		\eta(D^2 u)(0) = \begin{pmatrix}
			u_{nn} I_{n-1} + S & \xi \\
			\xi^T & \zeta
		\end{pmatrix},
		\qquad |S| + |\xi| + |\zeta| \le C.
		\]
		Admissibility implies $D^2 u(0) \in \widetilde{\Gamma}_1$, so $\Delta u(0) > 0$. Since tangential derivatives are bounded, this gives $u_{nn} \ge -C$.
		
		For the quotient equation, $\sigma_k(\eta) = \sigma_l(\eta)$. For large $u_{nn}$, the leading terms give
		\[
		C_{n-1}^k  u_{nn}^k - C u_{nn}^{k-1} \le C_{n-1}^l  u_{nn}^l + C u_{nn}^{l-1}.
		\]
		Since $k > l$, this yields $u_{nn} \le C$. Hence $|u_{nn}(0)| \le C$. Therefore, $|D^2 u(0)| \le C$.
		
		For $k=n$, set
		\[
		T := \sum_{\alpha < n} u_{\alpha\alpha}(0).
		\]
		Since $\Phi_n(0) = 0$, summing  over $\alpha = 1, \dots, n-1$ gives
		\[
		T = \sum_{\alpha < n} \Phi_{\alpha\alpha}(0) - u_n(0) \sum_{\alpha < n} \rho_{\alpha\alpha}(0).
		\]
		The quantity $-\sum_{\alpha < n} \rho_{\alpha\alpha}(0)$ is the mean curvature of $\partial D$ at $p$ with respect to the outward unit normal of $D$. By strict mean convexity and the lower bound $D_\nu u_R \ge K$ , we obtain
		\[
		0 < c_3 \le T \le C, \qquad c_3 = c_3(\partial D, K) > 0.
		\]
		Also,
		\[
		\eta(D^2 u)(0) = \begin{pmatrix}
			u_{nn} I_{n-1} + S & \xi \\
			\xi^T & T
		\end{pmatrix},
		\qquad |S| + |\xi|+|T| \le C.
		\]
		Admissibility gives $\eta(D^2 u)(0) > 0$, so in particular $u_{nn} \ge -C$. Since $F(D^2 u) = 1$ and $k=n$, we have $\sigma_n(\eta(D^2 u)(0)) = \sigma_l(\eta(D^2 u)(0))$, i.e.
		\[
		\det(\eta(D^2 u)(0)) = \sigma_l(\eta(D^2 u)(0)).
		\]
		\noindent\textbf{Case 1: $l \le n-2$.}
		For the quotient equation, $\sigma_n(\eta) = \sigma_l(\eta)$ with $k=n$.
		For large $u_{nn}$, the block form gives
		\[
		\det(\eta) \ge T (u_{nn} - C)^{n-1},
		\]
		and
		\[
		\sigma_l(\eta) \le C_{n-1}^l u_{nn}^l + O(u_{nn}^{l-1}).
		\]
		Thus
		\[
		T (u_{nn} - C)^{n-1} \le C_{n-1}^l u_{nn}^l + O(u_{nn}^{l-1}).
		\]
		Since $l \le n-2$ and $T \ge c_3 > 0$, the left-hand side is of order $u_{nn}^{n-1}$ while the right-hand side is of order $u_{nn}^l$ with $l < n-1$. This is impossible for $u_{nn}$ sufficiently large. Hence $u_{nn} \le C$.
		
		
		
		\noindent\textbf{Case 2: $l=n-1$.}
		Using the block form above,
		\[
		\det(\eta(D^2 u)(0)) = T \det\left( u_{nn} I_{n-1} + S - T^{-1} \xi \xi^T \right).
		\]
		Since $T$ is bounded above and below by positive constants, and $S - T^{-1}\xi\xi^T$ is uniformly bounded, the left-hand side grows like
		Set $M:=S-T^{-1}\xi\xi^T$. Then
		\[
		\det(\eta)=T\det(u_{nn}I_{n-1}+M).
		\]
		Expanding the determinant,
		\begin{equation}
			\det(u_{nn} I_{n-1} + M)
			=
			u_{nn}^{n-1}
			+
			\sigma_1(M) u_{nn}^{n-2}
			+
			\sigma_2(M) u_{nn}^{n-3}
			+
			\cdots
			+
			\sigma_{n-1}(M).
		\end{equation}
		Hence
		\begin{equation}
			\det(\eta)
			=
			T u_{nn}^{n-1}
			+
			T\sigma_1(M) u_{nn}^{n-2}
			+
			T\sigma_2(M) u_{nn}^{n-3}
			+
			\cdots
			+
			T\sigma_{n-1}(M).
		\end{equation}
		On the other hand,
		\begin{equation}
			\sigma_{n-1}(\eta)
			=
			u_{nn}^{n-1}
			+
			\big((n-1)T+\operatorname{tr}(S)\big)u_{nn}^{n-2}
			+
			O(u_{nn}^{n-3}).
		\end{equation}
		Suppose, for contradiction, that $u_{nn}$ is unbounded on $\partial D$. 
		Then there exists a sequence of points $p_j \in \partial D$ such that
		\begin{equation}
			u_{nn}(p_j) \to \infty.
		\end{equation}
		By compactness of $\partial D$, passing to a subsequence if necessary, 
		we may assume $p_j \to p_0$ for some $p_0 \in \partial D$.
		
		For each $j$, the block decomposition at $p_j$ gives
		\[
		\eta(D^2u)(p_j)=\begin{pmatrix}
			u_{nn}(p_j)I_{n-1}+S_j & \xi_j\\
			\xi_j^T & T_j
		\end{pmatrix},
		\]
		where $T_j:=\sum_{\alpha<n}u_{\alpha\alpha}(p_j)$ satisfies $0<c_3\le T_j\le C$, and $|S_j|+|\xi_j|\le C$ by the tangential derivative estimates.
		
		From the expansions
		\begin{equation}
			\det(\eta_j)
			=
			T_j u_{nn}^{n-1}(p_j)
			+
			\bigl(T_j\operatorname{tr}(S_j)-|\xi_j|^2\bigr)u_{nn}^{n-2}(p_j)
			+
			O(u_{nn}^{n-3}(p_j)).
		\end{equation}
		and
		\begin{equation}
			\sigma_{n-1}(\eta_j)
			=
			u_{nn}^{n-1}(p_j)
			+
			\bigl((n-1)T_j+\operatorname{tr}(S_j)\bigr)u_{nn}^{n-2}(p_j)
			+
			O(u_{nn}^{n-3}(p_j)).
		\end{equation}
		together with $\det(\eta_j)=\sigma_{n-1}(\eta_j)$, comparing the $u_{nn}^{n-1}(p_j)$-terms yields $T_j\to1$. Then comparing the $u_{nn}^{n-2}(p_j)$-terms yields
		\[
		\operatorname{tr}(S_j)-|\xi_j|^2 = (n-1)T_j+\operatorname{tr}(S_j)+o(1).
		\]
		Since $T_j\to1$ and $S_j,\xi_j$ are bounded, passing to a subsequence if necessary, we may assume $S_j\to S_0$ and $\xi_j\to\xi_0$. Taking the limit gives
		\[
		\operatorname{tr}(S_0)-|\xi_0|^2 = (n-1)+\operatorname{tr}(S_0),
		\]
		which implies
		\begin{equation}
			-|\xi_0|^2 = n-1.
		\end{equation}
		This is impossible since the left-hand side is non-positive while the right-hand side is positive for $n\ge3$. Hence $u_{nn}$ cannot tend to infinity, so $u_{nn}\le C$ on $\partial D$. Therefore $|D^2u_R|\le C$ on $\partial D$.
		\paragraph{Outer boundary $\partial E_R$.}
		We rescale the outer boundary to a fixed ellipsoid. Set
		\[
		\widehat{E} := \{ y \in \mathbb{R}^n : s(y) < 1 \}, \qquad y := R^{-1/2} x,
		\qquad \widetilde{u}_R(y) := R^{-1} u_R(\sqrt{R} y), \qquad \widehat{\Omega}_R := \widehat{E} \setminus R^{-1/2} \overline{D}.
		\]
		Then $D_y^2 \widetilde{u}_R(y) = D_x^2 u_R(\sqrt{R} y)$, so $\widetilde{u}_R$ satisfies the same quotient equation
		\[
		F(D_y^2 \widetilde{u}_R) = 1 \quad \text{in } \widehat{\Omega}_R.
		\]
		The rescaled comparison functions satisfy $\widetilde{q}_R \le \widetilde{u}_R \le \widetilde{\bar{u}}_R$ and all three are constant on $\partial \widehat{E}$. By the same argument as on the inner boundary, we obtain
		$|D^2 \widetilde{u}_R|\le C$ on $ \partial \widehat{E}.$
		Scaling back,we obtain
		\begin{equation}
			|D^2 u_R| \le C \quad \text{on } \partial E_R.
		\end{equation}
	\end{proof}

\end{document}
